% 三层物理架构
\section{三层物理架构}\label{sec:bg-arch}

晶圆级交换机把传统多机箱交换机的物理层级压缩进一块晶圆，其物理结构可自下而上分为三个层级。

\subsection{Die--Interposer 界面：$\mu$bump 键合与 die 间互联}\label{sec:bg-die-interposer}

每个功能裸片（die）通过面阵列 $\mu$bump 键合到硅 interposer 上，
pitch 可达 36--55\,$\mu$m，互联密度比 PCB 高出两个数量级~\cite{ucie2.0-2024}。
裸片之间的通信在 interposer 金属层上走线完成，对应 UCIe 链路~\cite{ucie2.0-2024}：
UCIe Advanced 面向 2.5D 短距低功耗场景，每 bit 功耗低至 0.25\,pJ。
这一层承载了交换拓扑中绝大多数组内链路，是超大 radix 的物理基础。

\subsection{Interposer--Substrate 界面：C4 焊球与出封装链路}\label{sec:bg-interposer-substrate}

interposer 经 C4 焊球连接到封装基板（substrate）。
需要跨组或出封装的链路（组间 SerDes 链路）在此经 C4 引出，再进入基板或背板走线，
对应 SerDes VSR/MR/LR 的中远距 reach 档位~\cite{oif-cei-5.1}。
C4 焊球数量有限，其中分配给 SerDes 信号的份额构成组间总出口带宽的硬约束。

\subsection{晶圆级集成尺度：垂直热网络}\label{sec:bg-wafer-scale}

在晶圆级尺度上，interposer 即整块晶圆，die 阵列分布其上；
热量沿 die $\to$ interposer $\to$ substrate $\to$ 环境的垂直路径逐层传导，
形成三层垂直热网络。组间 SerDes 的功耗抬升 substrate 温度，进而经 interposer 底面抬升组内 die 温度——
这一垂直耦合是第~\ref{chap:overview}章联合模型物理含义的来源。
